\documentclass[UTF8,fontset=none,aspectratio=169,10pt]{beamer}
\usepackage{ads-beamer-theme}
\title{高级分布式系统}
\subtitle{第十五讲：分布式共享内存与一致性模型}
\author{陈鹏飞\quad 计算机学院}
\institute{中山大学}
\date{2025--2026 学年}

\begin{document}
\begin{frame}[plain]
\begin{tikzpicture}[remember picture,overlay]
  \fill[ADSBlue] (current page.north west) rectangle ([xshift=2.4cm]current page.south west);
  \fill[ADSOrange] ([xshift=2.4cm]current page.north west) rectangle ([xshift=2.55cm]current page.south west);
  \node[anchor=west,align=left] at ([xshift=3.3cm,yshift=1.0cm]current page.west) {\fontsize{32}{38}\selectfont\color{ADSBlue}\bfseries 高级分布式系统};
  \node[anchor=west,align=left] at ([xshift=3.3cm,yshift=-0.4cm]current page.west) {\fontsize{21}{27}\selectfont\color{ADSDarkBlue}第十五讲\quad 分布式共享内存};
  \node[anchor=west,align=left] at ([xshift=3.3cm,yshift=-2.2cm]current page.west) {\large\color{ADSGray}把消息系统伪装成 Read/Write：方便，还是危险？};
  \node[anchor=east] at ([xshift=-0.8cm,yshift=-3.2cm]current page.east) {\small\color{ADSGray}陈鹏飞\quad 中山大学\quad 2025--2026};
  \draw[ADSBlue,line width=0.8pt] ([xshift=3.3cm,yshift=-1.35cm]current page.west) -- ([xshift=12.6cm,yshift=-1.35cm]current page.west);
\end{tikzpicture}
\end{frame}

\sectionframe{一、DSM 的承诺与代价}{应用看到一个共享地址空间，底层仍在做缓存、复制、消息传递和一致性维护。}

\begin{frame}{本讲学习目标}
\begin{enumerate}
  \item 理解 DSM 如何把消息传递隐藏在 Read/Write 接口后；
  \item 区分共享虚拟内存、页式 DSM、共享变量和共享对象；
  \item 掌握 coherence、linearizability、sequential consistency、causal consistency；
  \item 解释复制、缓存、失效和总序广播的代价；
  \item 把 DSM 思想映射到分布式缓存、参数服务器和 AI KV Cache。
\end{enumerate}
\vfill
\callout{主线}{接口越像本地内存，系统越需要在后台付出一致性和网络成本。}
\end{frame}

\begin{frame}{DSM 的抽象}
\begin{center}
\begin{tikzpicture}[node distance=10mm and 12mm,>=Latex]
  \nodebox[minimum width=2.4cm]{ADSBlue}{p1m}{Process 1};
  \nodebox[right=of p1m,minimum width=2.6cm]{ADSOrange}{m1}{Memory Manager};
  \nodebox[right=of m1,minimum width=2.6cm]{ADSGreen}{mem}{Shared Virtual Memory};
  \nodebox[below=of mem,minimum width=2.6cm]{ADSGray}{rep}{Replica/Cache};
  \draw[\arrow,ADSGray] (p1m)--node[above,font=\scriptsize]{Read/Write}(m1);
  \draw[\arrow,ADSGray] (m1)--(mem);
  \draw[\arrow,ADSGray] (mem)--(rep);
\end{tikzpicture}
\end{center}
\begin{itemize}
  \item 应用不直接调用 Send/Receive；
  \item DSM manager 在后台完成远程访问、复制和同步。
\end{itemize}
\end{frame}

\begin{frame}{DSM 的优点}
\begin{itemize}
  \item 单一地址空间，复杂对象可以按引用传递；
  \item 局部性好时，移动/缓存一整块数据可以降低远程访问次数；
  \item 屏蔽消息细节，编程接口更接近多线程共享内存；
  \item 可以使用普通网络和商用机器，扩大可用内存空间。
\end{itemize}
\end{frame}

\begin{frame}{DSM 的代价}
\begin{itemize}
  \item 应用必须理解一致性模型，否则“读到旧值”会造成逻辑错误；
  \item 底层仍然是异步消息传递，不会比手写的专用消息协议更快；
  \item 缓存失效、写入传播和网络分区使尾延迟不可预测；
  \item DSM manager 接管了通信路径，应用难以使用自定义优化。
\end{itemize}
\vfill
\callout{工程判断}{DSM 不是消除分布式复杂性，而是把复杂性从应用代码移到运行时和一致性契约。}
\end{frame}

\sectionframe{二、Coherence 与一致性模型}{同一个变量被多个副本读取和写入时，“最近一次写入”需要明确语义。}

\begin{frame}{为什么存在多个可能交错？}
\begin{itemize}
  \item P1 有 $s_1$ 个操作，P2 有 $s_2$ 个操作，……；
  \item 全局交错数量为 $(s_1+\cdots+s_n)!/(s_1!\cdots s_n!)$；
  \item 一致性模型定义哪些交错对应用可见；
  \item 模型越强，可见交错越少，协调与通信成本越高。
\end{itemize}
\end{frame}

\begin{frame}{Linearizability：像一个全局原子时钟}
\begin{itemize}
  \item 每个 Read/Write 像在某个瞬间原子发生；
  \item 不重叠操作必须保持真实时间顺序；
  \item 所有进程看到同一个全序；
  \item 适合锁、计数器、主备切换等需要强语义的对象。
\end{itemize}
\vfill
\callout{代价}{模拟全局时间轴通常需要总序广播、主节点或强同步。}
\end{frame}

\begin{frame}{Sequential Consistency}
\begin{itemize}
  \item 所有操作存在一个全局顺序；
  \item 每个进程的程序顺序在全局顺序中保持；
  \item 但不要求全局顺序服从真实墙上时间；
  \item 因此比 linearizability 弱，却通常更易实现。
\end{itemize}
\end{frame}

\begin{frame}{一个读写反例}
初始 $x=y=0$：
\begin{itemize}
  \item P1：Write$(x,4)$，随后 Read$(y)$ 返回 2；
  \item P2：Write$(y,2)$，随后 Read$(x)$ 返回 0；
  \item 这种结果不能由一个保持各自程序顺序的全局交错解释，因此不满足 sequential consistency；
  \item 若 P2 的 Read$(x)$ 返回 4，则可以构造一个一致全局顺序。
\end{itemize}
\end{frame}

\begin{frame}{Causal Consistency}
\begin{itemize}
  \item 只要求因果相关的写入在所有副本上顺序一致；
  \item 并发写入可以在不同副本看到不同顺序；
  \item 适合社交动态、协作编辑和跨区域缓存；
  \item 因果元数据可以用向量时钟、版本向量或逻辑依赖摘要表示。
\end{itemize}
\vfill
\callout{理解}{“更弱”不是“不可靠”，而是把不必要的全局顺序交给应用自行处理。}
\end{frame}

\sectionframe{三、复制与缓存：读快了，写复杂了}{DSM 的性能来自本地读；一致性成本主要发生在写入、失效和冲突解决。}

\begin{frame}{全复制与部分复制}
\begin{tabularx}{\linewidth}{>{\bfseries}p{2.8cm}XX}
\toprule
全复制 & 每个节点保存变量副本；读快，写需要传播到所有节点。 \\
部分复制 & 只有部分节点保存数据；降低容量/传播成本，但远程读更复杂。 \\
只读复制 & 适合模型、配置和静态数据；写入通过新版本发布。 \\
\bottomrule
\end{tabularx}
\end{frame}

\begin{frame}{写失效协议的基本流程}
\begin{enumerate}
  \item P1 写入 x，先获得写权限或版本号；
  \item DSM manager 向持有旧副本的节点发送 invalidate/update；
  \item 节点确认已失效或应用新版本；
  \item P1 完成写入并返回；
  \item 后续读请求从新版本或指定 leader 获取。
\end{enumerate}
\end{frame}

\begin{frame}{总序广播实现线性化}
\begin{itemize}
  \item Memory Manager 把 Read/Write 请求总序广播给所有副本；
  \item 所有副本按同样顺序执行 Write；
  \item Read 也参与总序，即使其他副本上是 no-op；
  \item 这样 Read 与并发 Write 的相对顺序不会分叉。
\end{itemize}
\vfill
\callout{连接第十四讲}{总序广播通常由共识日志或原子广播实现；DSM 的强一致性直接依赖共识。}
\end{frame}

\begin{frame}{为什么读操作也要排序？}
\begin{itemize}
  \item 若 Read 不进入总序，P1 可能在 Write(x,4) 之后读到 4；
  \item P2 同时在旧副本读到 0；
  \item 即使所有 Write 最终传播，两个 Read 的观察仍可能无法放入同一个线性化顺序；
  \item 将 Read 纳入总序可以定义其线性化点。
\end{itemize}
\end{frame}

\sectionframe{四、DSM 与现代分布式数据系统}{今天的缓存、参数服务器和模型服务，本质上都在选择“共享状态 + 一致性契约”。}

\begin{frame}{DSM 与分布式缓存}
\begin{itemize}
  \item Redis/Memcached 把共享状态暴露成键值 Read/Write；
  \item cache invalidation、TTL、版本号和租约对应 DSM 的复制一致性问题；
  \item 强一致缓存需要 leader/quorum，弱一致缓存允许短暂旧读；
  \item 应用要明确：读旧值会影响安全、计费还是只影响展示。
\end{itemize}
\end{frame}

\begin{frame}{参数服务器与分布式训练}
\begin{itemize}
  \item 参数服务器/embedding 服务提供共享参数空间；
  \item 同步训练追求每个 step 的 barrier 和一致版本；
  \item 异步训练允许 stale gradients，以换取吞吐和容错；
  \item 这正是强一致 DSM 与因果/最终一致缓存之间的工程选择。
\end{itemize}
\end{frame}

\begin{frame}{LLM KV Cache 的一致性边界}
\begin{itemize}
  \item KV Cache 通常绑定请求、模型版本和 GPU 实例；
  \item PD 分离时由 Prefill 产生、Decode 接管 KV；共享缓存需要会话亲和、版本标记和淘汰策略；
  \item 迁移到另一张 GPU 时，缓存可能不可直接复用；
  \item 让缓存“最终一致”通常可接受，但模型版本和安全权限必须强一致。
\end{itemize}
\end{frame}

\begin{frame}{近期案例：Flink 的 checkpoint barrier 与 watermark}
\begin{itemize}
  \item checkpoint barrier 沿数据流传播，用来把多个输入分区的状态切到同一个一致恢复边界；
  \item watermark 表示事件时间处理进度，用来判断窗口何时可以关闭；它是进度估计，不是 checkpoint；
  \item 背压或乱序会让 barrier、watermark 和数据记录以不同速度到达，运行时必须分别处理对齐和迟到数据；
  \item 这说明“状态一致性”和“业务时间进度”是两个不同的契约，不能用一个时间戳替代。
\end{itemize}
\vfill
\callout{课堂问题}{一个窗口已经收到 watermark，但对应 checkpoint 尚未提交：可以对外发布结果吗？需要哪些幂等或回滚条件？}
\end{frame}

\begin{frame}{课堂互动：为三种数据选择一致性}
\begin{tabularx}{\linewidth}{>{\bfseries}p{3.1cm}XX}
\toprule
模型生产 alias & 线性化/共识日志 & 不能出现两个生产版本 \\
推荐流缓存 & 因果/最终一致 & 允许短暂旧数据，追求低延迟 \\
训练 optimizer state & step 一致/Checkpoint barrier & 恢复后必须与模型参数匹配 \\
\bottomrule
\end{tabularx}
\vfill
\callout{答题要求}{先写出“如果读到旧值会造成什么后果”，再选择一致性强度。}
\end{frame}

\sectionframe{五、DSM 完整推演}{从一次旧读开始，追踪副本、失效、版本和一致性承诺。}

\begin{frame}{完整例子：库存扣减的两个副本}
初始 $stock=1$，副本 $R_A$、$R_B$ 都为 1：
\begin{center}
\begin{tikzpicture}[>=Latex,node distance=18mm]
\nodebox[minimum width=2.5cm]{ADSBlue}{ra}{副本 $R_A$\\stock=1};
\nodebox[right=of ra,minimum width=2.5cm]{ADSOrange}{rb}{副本 $R_B$\\stock=1};
\nodebox[below=of ra,minimum width=2.5cm]{ADSGreen}{ua}{用户 A};
\nodebox[below=of rb,minimum width=2.5cm]{ADSGray}{ub}{用户 B};
\draw[\arrow,ADSGray] (ua)--node[left,font=\scriptsize]{Read=1, Write=0}(ra);
\draw[\arrow,ADSGray] (ub)--node[right,font=\scriptsize]{Read=1, Write=0}(rb);
\draw[ADSOrange,dashed,->] (ra)--node[above,font=\scriptsize]{失效/更新}(rb);
\end{tikzpicture}
\end{center}
\begin{itemize}
  \item 两个用户都读到 1，再各自写入 0；
  \item 若没有原子写权限或冲突检测，系统会“成功扣减两次”，但库存只剩 0。
\end{itemize}
\end{frame}

\begin{frame}{这个执行满足哪种一致性？}
\begin{itemize}
  \item 若两个写入都成功，无法放入一个保持每个进程程序顺序的全局原子顺序；
  \item 线性化对象应拒绝其中一次写，或返回冲突；
  \item 最终一致副本可以暂时接受不同值，但必须提供冲突解决规则；
  \item “读到旧值”不是抽象问题：它会直接改变计费、库存和权限结果。
\end{itemize}
\vfill
\callout{判断步骤}{先写操作历史，再问是否存在一个合法全局顺序；不要只看最终两个副本是否相同。}
\end{frame}

\begin{frame}{写失效协议的逐步时间线}
\begin{enumerate}
  \item P1 向 manager 请求 $x$ 的写租约或新版本号 8；
  \item manager 向持有旧副本的 P2、P3 发送 invalidate(7)；
  \item P2、P3 确认失效后，P1 提交 Write$(x,8)$；
  \item 后续 Read 必须从版本 $\geq 8$ 的副本返回；
  \item P2 若再次读取，先回源或接收 update(8)。
\end{enumerate}
\vfill
\smallnote{写失效降低了传播 payload，但增加了确认路径；写更新则相反。}
\end{frame}

\begin{frame}{版本号如何防止旧副本复活？}
\begin{itemize}
  \item 每次写入携带单调递增版本或 epoch；
  \item 副本只接受版本更大的更新；
  \item 节点从网络分区恢复时，旧版本不能覆盖新版本；
  \item 若出现并发同版本写，需要 tie-breaker、向量版本或应用冲突合并。
\end{itemize}
\begin{center}
\begin{tikzpicture}[>=Latex,node distance=12mm]
\nodebox{ADSGray}{old}{旧副本 v7};
\nodebox[right=of old]{ADSBlue}{new}{新副本 v8};
\draw[\arrow,ADSOrange] (old)--node[above,font=\scriptsize]{拒绝覆盖}(new);
\end{tikzpicture}
\end{center}
\end{frame}

\begin{frame}{线性化与顺序一致性的具体差别}
\begin{enumerate}
  \item P1 在真实时间先完成 Write$(x,1)$；
  \item 随后 P2 发起 Read$(x)$，却返回 0；
  \item 若两操作不重叠，这违反 linearizability 的真实时间顺序；
  \item 若操作重叠，sequential consistency 仍可能允许把 Read 放在 Write 之前；
  \item 所以“全局顺序”不等于“墙上时间顺序”。
\end{enumerate}
\end{frame}

\begin{frame}{参数服务器：同步 step 的一致性边界}
\begin{center}
\begin{tikzpicture}[>=Latex,node distance=12mm and 15mm]
\nodebox[minimum width=2.2cm]{ADSBlue}{w1p}{Worker 1\\grad$_1$};
\nodebox[right=of w1p,minimum width=2.2cm]{ADSOrange}{ps}{Parameter\\Server v10};
\nodebox[right=of ps,minimum width=2.2cm]{ADSGreen}{w2p}{Worker 2\\grad$_2$};
\draw[\arrow,ADSGray] (w1p)--node[above,font=\scriptsize]{push}(ps);
\draw[\arrow,ADSGray] (w2p)--node[above,font=\scriptsize]{push}(ps);
\draw[\arrow,ADSBlue] (ps)--node[below,font=\scriptsize]{pull v11}(w1p);
\draw[\arrow,ADSOrange] (ps)--node[below,font=\scriptsize]{pull v11}(w2p);
\end{tikzpicture}
\end{center}
\begin{itemize}
  \item 同步训练要求一个 step 的 worker 使用同一版本参数；
  \item 异步训练允许 stale gradient，吞吐更高但收敛分析更复杂；
  \item 版本号和 step barrier 是 DSM 一致性契约的工程化实现。
\end{itemize}
\end{frame}

\begin{frame}{LLM KV Cache 的一次迁移}
\begin{enumerate}
  \item 请求在 GPU-A 上生成前缀缓存，key 包含模型版本 v3；
  \item 调度器把请求迁移到 GPU-B；
  \item GPU-B 发现本地没有 v3 缓存，从共享缓存读取或重新计算；
  \item 模型升级到 v4 后，v3 缓存不能被误认为 v4 命中；
  \item 租户 ID、会话 ID 和权限必须同时绑定在 cache key 上。
\end{enumerate}
\end{frame}

\begin{frame}{课堂练习：选择一致性并说明后果}
\begin{tabularx}{\linewidth}{>{\bfseries}p{3.2cm}XX}
\toprule
对象 & 可接受语义 & 读到旧值的后果 \\
生产模型 alias & 线性化/共识 & 两个版本同时对外服务 \\
推荐流缓存 & 因果/最终一致 & 页面短暂显示旧推荐 \\
optimizer state & step barrier & 恢复后梯度与参数错位 \\
KV cache & 版本一致 + 会话亲和 & 命中错误模型或跨租户泄露 \\
\bottomrule
\end{tabularx}
\end{frame}

\begin{frame}{DSM 调试清单}
\begin{enumerate}
  \item 记录每次 Read/Write 的副本、版本、请求 ID 和线性化候选点；
  \item 记录失效、租约、回源和冲突解决事件；
  \item 用同一 trace 重放强一致和弱一致配置；
  \item 对 P99 延迟分解：本地命中、远程读、失效等待、重试；
  \item 把“旧读是否允许”写进 API 契约和测试断言。
\end{enumerate}
\end{frame}

\sectionframe{六、一致性选择实训}{把“强一点还是弱一点”转化为可测量的读写、修复和成本决策。}

\begin{frame}{Quorum 读写的基本条件}
设副本数为 $N$，读 quorum 为 $R$，写 quorum 为 $W$：
\begin{itemize}
  \item 若 $R+W>N$，一次读集合与最近一次写集合必相交；
  \item 若 $W>N/2$，两个写 quorum 也必相交；
  \item quorum 只提供交集条件，不自动解决网络分区、冲突合并和版本选择；
  \item 读修复可以把旧副本逐步更新到最新版本。
\end{itemize}
\end{frame}

\begin{frame}{读修复（read repair）例子}
\begin{enumerate}
  \item 客户端向 R1、R2、R3 读取 key=x；
  \item 返回版本分别为 v8、v8、v7；
  \item 客户端选择 v8 返回给应用；
  \item 后台向 R3 写回 v8；
  \item 下一次读取命中三个相同版本，降低继续读到旧值的概率。
\end{enumerate}
\vfill
\smallnote{读修复把一致性成本摊到读路径；若读流量很低，旧副本可能长期存在。}
\end{frame}

\begin{frame}{冲突合并：LWW 与 CRDT}
\begin{tabularx}{\linewidth}{>{\bfseries}p{3cm}XX}
\toprule
LWW & 按时间戳选最后写入。 & 简单，但时钟偏差可能丢更新。 \\
版本向量 & 保留并发版本，交给应用合并。 & 语义准确，元数据更大。 \\
CRDT & 设计可交换、幂等、结合的合并操作。 & 适合计数器/集合，不适合所有对象。 \\
\bottomrule
\end{tabularx}
\end{frame}

\begin{frame}{线性化历史练习}
初始 $x=0$：
\begin{enumerate}
  \item P1：Write$(x,1)$ 完成；
  \item P2：Read$(x)$ 返回 0；
  \item 若 P2 的调用发生在 P1 返回之后，这个历史违反 linearizability；
  \item 若两操作重叠，则需要进一步检查是否能构造合法全序。
\end{enumerate}
\end{frame}

\begin{frame}{最终一致缓存的读后果}
\begin{itemize}
  \item 页面推荐读到旧值：通常只影响体验；
  \item 权限配置读到旧值：可能扩大访问范围；
  \item 计费额度读到旧值：可能造成超额消费；
  \item 模型 alias 读到旧值：可能让流量同时落到两个生产版本；
  \item 先按后果分类，再决定是否需要强一致路径。
\end{itemize}
\end{frame}

\begin{frame}{参数服务器中的 stale gradient}
\begin{enumerate}
  \item Worker A 拉到参数版本 10，计算梯度耗时 5 秒；
  \item 期间 Worker B 已提交版本 11；
  \item A 将基于 v10 的梯度提交给服务器；
  \item 同步方案拒绝或重算，异步方案可能接受并记录 staleness=1；
  \item staleness 上界是训练稳定性与吞吐之间的可调参数。
\end{enumerate}
\end{frame}

\begin{frame}{Cache stampede：一致性与负载的耦合}
\begin{itemize}
  \item 热点 key 过期时，大量 worker 同时发现 miss；
  \item 每个 worker 都回源加载模型或 KV Cache，造成放大流量；
  \item single-flight、租约和提前刷新减少重复加载；
  \item 这些机制本质上是在缓存对象上增加互斥和版本一致性。
\end{itemize}
\end{frame}

\begin{frame}{课堂实验：同一 trace 的四种语义}
\begin{enumerate}
  \item 固定同一组 Read/Write 操作和网络延迟；
  \item 分别运行 linearizable、sequential、causal、eventual 配置；
  \item 记录旧读比例、冲突数、修复流量、P99 延迟；
  \item 解释哪一种语义刚好满足业务契约，哪一种是过度设计。
\end{enumerate}
\end{frame}

\begin{frame}{DSM 设计评审清单}
\begin{enumerate}
  \item 对象的 owner、版本和副本集合是什么？
  \item 读操作是否允许旧值？允许多久？
  \item 写冲突如何检测、合并或拒绝？
  \item 节点分区、恢复和缓存淘汰如何改变语义？
  \item 观测指标能否区分网络延迟、失效等待和回源重算？
\end{enumerate}
\end{frame}

\begin{frame}{本讲总结}
\begin{enumerate}
  \item DSM 用 Read/Write 隐藏消息传递，但一致性成本仍然存在；
  \item Linearizability 最强，Sequential/Causal/最终一致性逐步放宽全局顺序；
  \item 复制、缓存、失效和总序广播共同决定读写性能；
  \item DSM 的思想今天仍出现在分布式缓存、参数服务器、KV Cache 和模型注册中心；
  \item 选择一致性模型前必须先明确业务能接受什么旧读、冲突和重试。
\end{enumerate}
\vfill
\callout{下一讲衔接}{下一讲进入认证、授权、机密性和 Byzantine 行为，讨论“谁有权写共享状态”。}
\end{frame}

\begin{frame}{课后题}
\begin{enumerate}
  \item 给出一个满足 sequential consistency 但不满足 linearizability 的读写执行；
  \item 设计总序广播 DSM 的 Read/Write 伪代码；
  \item 比较参数服务器同步与异步训练的一致性/吞吐取舍；
  \item 为 LLM KV Cache 设计版本号、会话亲和和失效策略。
\end{enumerate}
\end{frame}

\end{document}
